\documentclass[10pt]{beamer}

% Setting up the Beamer theme and colors
\usetheme{Copenhagen}
\usecolortheme{seahorse}
\setbeamertemplate{navigation symbols}{}
\setbeamertemplate{footline}[frame number]

% Including necessary packages
\usepackage[utf8]{inputenc}
\usepackage[T1]{fontenc}
\usepackage{lmodern}
\usepackage{amsmath, amssymb, amsfonts}
\usepackage{booktabs}
\usepackage{graphicx}
\usepackage{xcolor}
\usepackage{hyperref}
\usepackage{siunitx}
\usepackage{multirow}
\usepackage{tikz}

% Configuring hyperref
\hypersetup{
    colorlinks=true,
    linkcolor=blue,
    urlcolor=blue,
    citecolor=blue
}

% Title slide information
\title{A Step Toward Quantifying Independently Reproducible Machine Learning Research}
\subtitle{NeurIPS 2019}
\author{Edward Raff \\[0.5cm] Presented by: Yongcheng Huang (5560950), Group 2A}
\institute{Booz Allen Hamilton \& University of Maryland, Baltimore County}
\date{May 12, 2025}

\begin{document}

% Title slide
\begin{frame}
    \titlepage
\end{frame}

% Outline slide
\begin{frame}{Outline}
    \tableofcontents
\end{frame}

\section{Why Interesting?}

% Motivation slide
\begin{frame}{Why Should We Care?}
    \begin{itemize}
        \item Reproducibility is critical for trust and progress in AI/ML research.
        \item Concerns about a reproducibility crisis undermine scientific credibility.
        \item Conferences like NeurIPS push for code sharing, but independent reproducibility remains unquantified.
        \item \textbf{Goal}: Understand what makes ML papers independently reproducible without relying on authors' code.
    \end{itemize}
\end{frame}

\section{How Done Now?}

% Current Approaches slide
\begin{frame}{Current Approaches to Reproducibility}
    \begin{itemize}
        \item Conferences encourage or mandate code submission (e.g., NeurIPS, ICML).
        \item Authors often share code on platforms like GitHub.
        \item Reproducibility challenges (e.g., ICLR) involve re-running authors' code.
        \item \textbf{Limitation}: These approaches rely on authors' implementations, not independent replication.
    \end{itemize}
\end{frame}

\section{What is Missing, and So What?}

% Problem Statement slide
\begin{frame}{What is Missing?}
    \begin{itemize}
        \item \textbf{Problem}: Lack of systematic study on independent reproducibility (reproducing without authors' code).
        \item Current metrics focus on code availability, not clarity or detail in the paper itself.
        \item \textbf{Consequences}:
            \begin{itemize}
                \item Unclear what paper features enable successful reproduction.
                \item Researchers may waste time on poorly documented methods.
                \item Hinders scientific progress and trust in published results.
            \end{itemize}
    \end{itemize}
\end{frame}

\section{Proposed Approach}

% Study Overview slide
\begin{frame}{Proposed Approach: Quantifying Reproducibility}
    \begin{itemize}
        \item Analyze 255 ML papers (1984--2017) for independent reproducibility.
        \item Manually implement each paper’s method without using authors’ code.
        \item Record 26 paper features (e.g., readability, equations, pseudo-code).
        \item Use statistical analysis to identify features correlated with reproducibility.
        \item \textbf{Outcome}: 63.5\% (162/255) papers successfully reproduced.
    \end{itemize}
\end{frame}

% Methodology slide
\begin{frame}{Methodology}
    \begin{itemize}
        \item \textbf{Paper Selection}:
            \begin{itemize}
                \item New algorithm/method, implemented 2012--2017, no prior code exposure.
                \item Excluded 2018 papers to avoid time bias.
            \end{itemize}
        \item \textbf{Reproducibility Criteria}:
            \begin{itemize}
                \item $\geq 75\%$ claims confirmed with independent code.
                \item Flexible thresholds for order-of-magnitude or ranking claims.
            \end{itemize}
        \item \textbf{Features}: Objective (e.g., page count), mildly subjective (e.g., tables), subjective (e.g., readability).
    \end{itemize}
\end{frame}

\section{Experimental Questions}

% Experimental Questions slide
\begin{frame}{Experimental Questions}
    \begin{itemize}
        \item Which paper features predict successful independent reproduction?
        \item Does readability impact reproducibility more than technical details like equations?
        \item Are empirical papers easier to reproduce than theoretical ones?
        \item Does author engagement (e.g., replying to queries) improve reproducibility?
        \item \textbf{Setup}: Controlled analysis of 26 features across 255 papers, using statistical significance (p $\leq$ 0.05).
    \end{itemize}
\end{frame}

\section{Key Findings}

% Significant Features slide
\begin{frame}{Key Findings: Significant Features (p $\leq$ 0.05)}
    \begin{table}
        \centering
        \footnotesize
        \begin{tabular}{l c}
            \toprule
            Feature & p-value \\
            \midrule
            Readability & $9.68 \times 10^{-25}$ \\
            Rigor vs. Empirical & $1.55 \times 10^{-9}$ \\
            Authors Reply & $6.01 \times 10^{-8}$ \\
            Hyperparameters Specified & $8.45 \times 10^{-6}$ \\
            Compute Needed & $8.75 \times 10^{-5}$ \\
            Pseudo-Code & $2.31 \times 10^{-4}$ \\
            Primary Topic & $7.04 \times 10^{-4}$ \\
            Algorithm Difficulty & $2.94 \times 10^{-5}$ \\
            Number of Tables & 0.010 \\
            Number of Equations & 0.004 \\
            \bottomrule
        \end{tabular}
    \end{table}
\end{frame}

% Readability Impact slide
\begin{frame}{Readability: Strongest Predictor}
    \begin{itemize}
        \item \textbf{Excellent} readability (1 read needed): 100\% reproducible.
        \item Low readability papers shorter by 3.17--5.67 pages (p = 0.035).
        \item More equations per page reduce reproducibility (linked to readability).
        \item \textbf{Implication}: Clear communication is critical for reproducibility.
    \end{itemize}
\end{frame}

% Other Findings slide
\begin{frame}{Other Findings}
    \begin{itemize}
        \item \textbf{Empirical/Balanced} papers more reproducible than theoretical ones.
        \item \textbf{Pseudo-Code} helps, but step-code can reduce readability.
        \item \textbf{Tables} aid reproducibility; graphs/plots do not.
        \item \textbf{Author Reply}: 22/26 papers with replies reproduced (p = $6.01 \times 10^{-8}$).
        \item \textbf{Non-Significant}: Year published, code availability, venue type, appendix presence.
    \end{itemize}
\end{frame}

\section{Implications}

% Implications slide
\begin{frame}{Implications for ML Research}
    \begin{itemize}
        \item Increase page limits to allow detailed descriptions.
        \item Prioritize readability in peer review criteria.
        \item Encourage author engagement to clarify queries.
        \item Re-evaluate pseudo-code: High-level steps can confuse.
        \item Support empirical papers with practical details.
        \item Improve cluster computing frameworks for reproducibility.
    \end{itemize}
\end{frame}

\section{Limitations}

% Limitations slide
\begin{frame}{Limitations}
    \begin{itemize}
        \item Single reproducer: Results tied to author’s background (JSAT focus).
        \item Selection bias: Papers chosen based on interest, excluding prior code exposure.
        \item Subjective features: Readability, difficulty rely on author’s judgment.
        \item Limited historical records: Reliant on paper cataloging notes.
        \item No detailed failure reasons recorded (e.g., unclear notation).
    \end{itemize}
\end{frame}

\section{Conclusion}

% Conclusion slide
\begin{frame}{Conclusion}
    \begin{itemize}
        \item First study quantifying independent reproducibility in ML (63.5\% success rate).
        \item Key drivers: Readability, empirical focus, author engagement, hyperparameters.
        \item No evidence of a recent reproducibility crisis; issues persistent since 1984.
        \item Future work: Multi-reproducer studies, standardized features, survival analysis.
    \end{itemize}
\end{frame}

% Thank You slide
\begin{frame}{Thank You}
    \begin{center}
        \Large Questions? \\
        \vspace{0.5cm}
        Contact: \href{mailto:raff_edward@bah.com}{raff_edward@bah.com} \\
        Data: \url{https://github.com/EdwardRaff/Quantifying-Independently-Reproducible-ML}
    \end{center}
\end{frame}

\end{document}